% 2022 全国乙卷（文）第 7 题程序框图，内容依据原始 PDF 第 1 页。
% Source PDF: source/普通高考/2022/2022全国乙文(河南,山西,江西,安徽,陕西,甘肃,青海,内蒙古,黑龙江,吉林,.宁夏,新疆).pdf.
% Grid evidence: tmp/repaint_2022_national_b_liberal/grid/page_grid100.png and page_grid50.png.
% No-grid source crop: tmp/repaint_2022_national_b_liberal/crop/q7_source_tight.png.
% Node-edge list: 开始→输入→b=b+2a→a=b-a,n=n+1→判定;
% 判定“是”→输出n→结束；判定“否”→左侧回边→输入后、b更新前的主干汇合点。
% solid segments: all node boundaries and connector shafts; dashed segments: none.
% Labels are attached to their branch edges; every node clears problem paragraph state and is centered.

\begin{tikzpicture}[
  x=1cm,
  y=0.8cm,
  >=Stealth,
  line width=0.45pt,
  line cap=round,
  line join=round,
  every node/.style={font=\small,anchor=center,align=center,inner sep=0pt,
    execute at begin node={\setlength{\parindent}{0pt}\setlength{\leftskip}{0pt}\setlength{\rightskip}{0pt}\everypar{}}},
  flow/.style={draw,anchor=center,align=center,inner sep=0pt,
    execute at begin node={\setlength{\parindent}{0pt}\setlength{\leftskip}{0pt}\setlength{\rightskip}{0pt}\everypar{}}},
  startstop/.style={flow,rounded corners=0.18cm,minimum width=1.00cm,minimum height=0.54cm},
  process/.style={flow,minimum width=1.95cm,minimum height=0.56cm,inner xsep=4pt,inner ysep=2pt},
  inputoutput/.style={flow,shape=trapezium,trapezium left angle=72,trapezium right angle=108,
    minimum width=3.30cm,minimum height=0.58cm,inner xsep=4pt,inner ysep=2pt},
  decision/.style={flow,shape=diamond,aspect=2.8,minimum width=5.35cm,minimum height=1.34cm,
    inner xsep=4pt,inner ysep=2pt},
  branchlabel/.style={anchor=center,align=center,inner sep=0pt,
    execute at begin node={\setlength{\parindent}{0pt}\setlength{\leftskip}{0pt}\setlength{\rightskip}{0pt}\everypar{}}},
]
  \node[startstop] (start) at (0,8.95) {开始};
  \node[inputoutput] (input) at (0,7.72) {输入 $a=1,b=1,n=1$};
  \node[process] (updateb) at (0,6.03) {$b=b+2a$};
  \node[process,minimum width=2.90cm] (updatean) at (0,4.80) {$a=b-a,n=n+1$};
  \node[decision] (check) at (0,3.28) {$\left|\dfrac{b^2}{a^2}-2\right|<0.01$};
  \node[inputoutput,minimum width=1.80cm] (output) at (0,1.44) {输出 $n$};
  \node[startstop] (finish) at (0,0.14) {结束};

  \draw[->] (start.south) -- (input.north);
  \draw[->] (input.south) -- (updateb.north);
  \draw[->] (updateb.south) -- (updatean.north);
  \draw[->] (updatean.south) -- (check.north);
  \draw[->] (check.south) -- node[branchlabel,anchor=west,pos=0.18,xshift=0.08cm] {是} (output.north);
  \draw[->] (output.south) -- (finish.north);

  % 否分支左回，箭头在输入后、b 更新前的共享入口，和主干下行箭头独立。
  % The return arrow ends at the shared entrance; no plain shaft follows its tip.
  \coordinate (loopJoin) at (0,6.87);
  \coordinate (loopBottom) at (-3.15,3.72);
  \coordinate (loopTop) at (-3.15,6.87);
  \draw (check.west) -- node[branchlabel,below,pos=0.24] {否} (loopBottom)
    -- (loopTop);
  \draw[->] (loopTop) -- (loopJoin);

  \path[use as bounding box] (-3.45,-0.20) rectangle (3.45,9.28);
\end{tikzpicture}
